% ============================================================================
%  Appendices
% ============================================================================
\appendix

\section{Test Vectors}
\label{app:test-vectors}

These test vectors are normative. A conformant implementation MUST decode the encoded forms and recover the listed binary forms; conversely, encoding the listed binary forms MUST produce the listed encoded forms.

\subsection*{Vector 1: Standard variant}

Inputs:
\begin{itemize}
  \item Timestamp: \code{1700000000000} (2023-11-14T22:13:20Z), hex \code{0x018BCFE56800}
  \item Variant: \code{0x01}
  \item Counter: \code{0x0042}
  \item Random tail: 32 bytes of \code{0xAA} (synthetic, for test purposes)
\end{itemize}

Binary form (41 bytes, hex):
\begin{lstlisting}[basicstyle=\ttfamily\scriptsize,frame=single]
01 8B CF E5 68 00  01  00 42  AA AA AA AA AA AA AA AA
                                AA AA AA AA AA AA AA AA
                                AA AA AA AA AA AA AA AA
                                AA AA AA AA AA AA AA AA
\end{lstlisting}

Base64URL: \code{AYvP5WgAAQBCqqqqqqqqqqqqqqqqqqqqqqqqqqqqqqqqqqqqqqqqqqo}

Crockford Base32: \code{065WZSB8000G0GNANANANANANANANANANANANANANANANANANANANANANANANANAN8}

\subsection*{Vector 2: HLC variant}

Inputs:
\begin{itemize}
  \item Logical timestamp: \code{1700000000123} (2023-11-14T22:13:20.123Z), hex \code{0x018BCFE5687B}
  \item Variant: \code{0x02}
  \item Counter: \code{0x0001}
  \item Random tail: 32 bytes of \code{0x55} (synthetic)
\end{itemize}

Binary form (41 bytes, hex):
\begin{lstlisting}[basicstyle=\ttfamily\scriptsize,frame=single]
01 8B CF E5 68 7B  02  00 01  55 55 55 55 55 55 55 55
                                55 55 55 55 55 55 55 55
                                55 55 55 55 55 55 55 55
                                55 55 55 55 55 55 55 55
\end{lstlisting}

Base64URL: \code{AYvP5Wh7AgABVVVVVVVVVVVVVVVVVVVVVVVVVVVVVVVVVVVVVVVVVVU}

Crockford Base32: \code{065WZSB8FC1000ANANANANANANANANANANANANANANANANANANANANANANANANANAM}

\section{Quick Reference Card}
\label{app:quickref}

\begin{center}
\renewcommand{\arraystretch}{1.2}
\begin{tabular}{ll}
\toprule
\textbf{Property} & \textbf{Value} \\
\midrule
Total size                  & 41 bytes \\
Timestamp width             & 48 bits (ms since Unix epoch) \\
Variant byte                & offset 6, \code{0x01} (Std) or \code{0x02} (HLC) \\
Counter width               & 16 bits, big-endian \\
Random tail width           & 256 bits (32 bytes) \\
Base64URL length            & 55 chars (no padding) \\
Crockford Base32 length     & 66 chars (no padding) \\
Time range                  & through year 10889 (8{,}919-year forward) \\
Classical PQ security       & 128 bits \\
Quantum (Grover) security   & 128 bits (second-preimage) \\
HLC drift bound             & 60{,}000 ms (normative max and default; tighten-only) \\
Max counter step            & 16 (random in [1,16]) \\
\bottomrule
\end{tabular}
\end{center}

\section{Example TesseraIDs (Same Millisecond)}
\label{app:examples}

The following 12 Crockford-encoded TesseraIDs (66 characters each) were generated within a single millisecond (timestamp \code{1781126753402}, 2026-06-10) by the reference Python implementation's \code{StandardGenerator}, called in a tight loop; the strings below are its verbatim \code{encode\_b32} output. They sort lexicographically in generation order, and each decodes to a valid 41-byte binary value via \code{decode\_b32}.

\begin{lstlisting}[basicstyle=\ttfamily\scriptsize]
06FB6VBCF80G0HZYGHBHEC3NZ8TGY0PZ6F6YR8THW28SHA3QANSCV9VX2C54EGSR1M
06FB6VBCF80G0MJ2K6VM5NFHTFEM4KGP5588C8BW0D1QSQV0WPZY4T5G88JT47Q5HC
06FB6VBCF80G0QTJZW65GE40SPEXBPER9RFZQZ49WGC2Y407NXHCCA49SB48MG5P30
06FB6VBCF80G0TPT7GWXHXQBTGB4R9WDNW4N6GVA2X7KGPJ2GW0KJ8182AH1AHB5PC
06FB6VBCF80G0V0D6137MM73659VF18B6N0SW2S89KVJM83NATAXFFZXBKGDSJS65M
06FB6VBCF80G0XR2JB2H8203XGKE8E75ND90ZTKDEPVD9S70BR76SRGSM755YWS3R4
06FB6VBCF80G0ZE7180J0QN0KPSXZ4F2RY4PQ8ACG13MDAW9J2EPBYWYD4V153PSDR
06FB6VBCF80G10H71NF10VJQFEKG83HYJK713YNYV6C4C5CJQGS6ET9A0RZ2G76KE4
06FB6VBCF80G12ETHBZQGP0SW9HEHY9R2P6KC1XV7D6W7EBXZT140NCS4SBZ21GYGC
06FB6VBCF80G13PY5Q61CXVC8T6PF6TMRMHSG0GZJYJCWPQJQPDX9DNSK0E7H7JS7G
06FB6VBCF80G148KMQKRGTZRHPZ6AYKADKM01WRYZC7KAKT8R9JVYNYJ4J8MNSK1QG
06FB6VBCF80G16X0AD6NQXCMS9JV978PPP8AND9KK28WMESJB50S0F0XMC05MW03EC
\end{lstlisting}

The first 12 characters (\code{06FB6VBCF80G}) are identical across all 12: those 60 bits cover the 48-bit timestamp, the \code{0x01} variant byte, and the zero high nibble of the counter. Only the remaining counter bits and the random tail differ. The shared prefix demonstrates the time-clustering property that makes TesseraID friendly to B-tree indexes.

\section{Glossary of Terms}
\label{app:glossary}

\textbf{Base32, Base64.} Binary-to-text encoding schemes. Base32 packs 5 bits per character (32-character alphabet); Base64 packs 6 bits per character (64-character alphabet). Defined in RFC 4648.

\textbf{Base64URL.} The URL-and-filename-safe variant of Base64 (RFC 4648 §5), using \code{A--Z}, \code{a--z}, \code{0--9}, \code{-}, \code{\_}. TesseraID's canonical text encoding.

\textbf{BHT (Brassard--H{\o}yer--Tapp).} A 1997 quantum collision-finding algorithm. Costs $\sim 2^{n/3}$ operations for an $n$-bit space, vs $\sim 2^{n/2}$ classically.

\textbf{Birthday bound.} The expected number of generations before a collision in a random sample. Classically $\sqrt{\pi/2} \cdot 2^{n/2}$ for an $n$-bit space.

\textbf{Crockford Base32.} Douglas Crockford's Base32 alphabet (\code{0--9} then \code{A--Z} minus \code{I/L/O/U}). Case-insensitive on decode, with I/L→1 and O→0 confusables canonicalized. ASCII-sort-preserving.

\textbf{CSPRNG.} Cryptographically Secure Pseudo-Random Number Generator. The OS source TesseraID generators draw from: \code{getrandom(2)} on Linux, \code{arc4random\_buf(3)} on macOS/BSD, \code{BCryptGenRandom} on Windows.

\textbf{Drift bound.} The maximum amount by which an HLC's logical time may exceed local physical time before generation is rejected. TesseraID: 60{,}000 ms, the normative maximum and default; per-generator overrides may only tighten it (§\ref{sec:spec}).

\textbf{Grover's algorithm.} A 1996 quantum search algorithm that reduces unstructured search from $O(N)$ to $O(\sqrt{N})$. Reduces second-preimage cost on an $n$-bit field from $2^n$ to $2^{n/2}$.

\textbf{HLC (Hybrid Logical Clock).} A clock combining physical time with a logical counter, preserving causal ordering across nodes that exchange clock values. TesseraID's variant 0x02. \citet{kulkarni2014hlc}.

\textbf{KSUID.} K-Sortable Globally Unique ID. 20-byte time-sortable identifier, 32-bit second timestamp + 128-bit random tail. Segment, 2017.

\textbf{ML-DSA, ML-KEM.} NIST's 2024 post-quantum cryptographic standards (FIPS 204 / FIPS 203). Reference points for TesseraID's 128-bit security target.

\textbf{Post-Quantum Security.} The cryptographic security level retained against quantum-equipped adversaries. For an $n$-bit random field: $n/2$ bits against second-preimage (Grover); $n/3$ against collision (BHT).

\textbf{Second Preimage.} The cost of finding a specific value in a random space (vs any colliding pair). The relevant threat model for identifier formats.

\textbf{Snowflake.} Twitter's 2010 64-bit identifier: 41-bit ms timestamp + 10-bit machine ID + 12-bit sequence. Requires coordinated machine-ID assignment.

\textbf{ULID.} Universally Unique Lexicographically Sortable Identifier. 16-byte time-sortable identifier, 48-bit ms timestamp + 80-bit random tail. Feerasta, 2016.

\textbf{UUIDv4.} Random Universally Unique Identifier (RFC 4122). 122 random bits, no ordering, no time component.

\textbf{UUIDv7.} Time-sortable UUID variant (RFC 9562, 2024). 48-bit ms timestamp + 74 random bits + version + variant marker. Sortable between milliseconds; intra-millisecond ordering is implementation-defined.

\textbf{Variant byte.} Byte at offset 6 of a TesseraID. \code{0x01} = Standard; \code{0x02} = HLC; other values reserved. The reserved-future range is shared across the \tessera{} monorepo: \code{0x03} is assigned to \tesserahash{}, \code{0x04} to \tesseraseal{}; \code{0x00} and \code{0xFF} are permanently reserved. TesseraID decoders MUST reject every byte they do not own.

\section{The \texorpdfstring{\tessera{}}{tessera} Monorepo Context}
\label{sec:appendix-monorepo}

TesseraID v1 is the first of three projects under the \tessera{} monorepo. Each project ships an independent normative wire format, its own whitepaper, and byte-identical polyglot reference implementations across the same eight languages (Python, C, C++, Go, Rust, TypeScript, Java, CPython extension). The projects share a single variant-byte namespace and a single cross-language conformance harness.

\subsection*{Project relationship}

\begin{center}
\small
\begin{tabular}{@{}lllr@{}}
\toprule
\textbf{Project} & \textbf{Concern} & \textbf{Variant byte} & \textbf{Wire size (bytes)} \\
\midrule
\tesseraid{}    & Post-quantum-era identifier            & \code{0x01} (Standard), \code{0x02} (HLC) & 41 \\
\tesserahash{}  & SHAKE-256 content hash                  & \code{0x03}                              & 73 \\
\tesseraseal{} & ML-KEM-1024 encryption envelope         & \code{0x04}, suite byte \code{0x01}/\code{0x02}/\code{0x06} & variable (header 16 + body) \\
\bottomrule
\end{tabular}
\end{center}

The first nine bytes of every project's wire format are layout-compatible: \code{TS(6) || VARIANT(1) || CTR(2)}. A single shared header \code{common/c/tessera\_header.h} pins those offsets across all three. Decoders in each project reject every variant byte they do not own; the rejection contracts are covered by \texttt{invalid.json} KAT files in every project's \texttt{conformance/vectors/}.

\subsection*{The ``\tesseraid{}-as-AD'' canonical pattern}

When a TesseraID identifies a record whose content needs confidentiality and integrity:

\begin{lstlisting}[language=Pseudocode]
record_id  = tessera_id.new_hlc()                              # 41 bytes
envelope   = tessera_seal.encrypt(plaintext, recipient_pub,
                             associated_data=record_id)  # binds the TesseraID
# ... after delivery ...
plaintext  = tessera_seal.decrypt(envelope, recipient_priv,
                             expected_ad=record_id)      # MUST match
\end{lstlisting}

Substituting the envelope under a different \tesseraid{} fails decryption (AEAD AD mismatch). This pattern costs zero additional cryptography --- the AEAD's AD path is free --- and binds the identifier and ciphertext into a single cryptographic guarantee. See the \tesseraseal{} whitepaper for the normative description.

When a TesseraID identifies a record whose \emph{content bytes} need an integrity check independent of any encryption, use \tesserahash{}: the \tesserahash{} of the plaintext can be carried alongside the TesseraID (or used as the AEAD AD in a \tesseraseal{} envelope) to bind the envelope to ``this specific content,'' not just ``some content this recipient can decrypt.''

\subsection*{Shared discipline}

All three projects share:

\begin{itemize}
\item One specification-first methodology: whitepaper LaTeX is normative, Python is the readable reference, C pins the canonical ABI, every other language ports byte-identical.
\item One cross-language KAT harness at \texttt{tools/conformance/run.sh}; CI's branch-protection gate.
\item One shared library directory at \texttt{common/} for primitives used by more than one project (Keccak/SHAKE-256 for C/C++/Java; HKDF-SHA3 and shared header offsets for \tesseraseal{}-adjacent code).
\item One author and one set of memory invariants: see Appendix~\ref{sec:appendix-provenance}.
\end{itemize}

\section{Author and Document Provenance}
\label{sec:appendix-provenance}

\textbf{Author:} Bahram Samani, \tessera{} monorepo maintainer.

This document is built from LaTeX sources under \texttt{tessera-id/latex/} (\texttt{main.tex} + \texttt{secNN\_*.tex}); the build chain is \texttt{pdflatex $\rightarrow$ bibtex $\rightarrow$ pdflatex $\times$ 2} per the Makefile. The artifact at the repo root \texttt{TesseraID\_Whitepaper.pdf} is regenerated and committed alongside any normative change.

\textbf{Authority order} (when sources disagree): this whitepaper, then the Python reference at \texttt{tessera-id/python/tessera\_id.py} (with its optimized sibling \texttt{tessera\_id\_fast.py}), then the canonical C ABI at \texttt{tessera-id/c/tessera\_id.h}, then every other language. A cross-language disagreement is never resolved by changing the spec to match an outlier; the outlier is fixed.

\textbf{Hard invariants} (from \texttt{tessera-id/CLAUDE.md}; security-relevant changes must update both):
\begin{itemize}
\item 41-byte layout: \code{TS(6) || VARIANT(1) || CTR(2) || RANDOM(32)}, big-endian.
\item Variant byte: \code{0x01} = Standard, \code{0x02} = HLC; everything else (including the cross-project \code{0x03}/\code{0x04}) rejected.
\item Canonical encodings: Base64URL (55 chars, no padding) and Crockford Base32 (66 chars, case-insensitive on decode). The final Crockford character carries 3 bits of value plus 2 bits of zero padding; Crockford decode rejects inputs whose 2 trailing low padding bits of that final character are non-zero (the implementations reject a non-zero low-2-bit pad).
\item HLC drift bound: \texttt{HLC\_MAX\_DRIFT\_MS = 60\_000} ms. Peer HLC more than 60 s ahead of local physical time is a hard \texttt{HLCDriftError}, not a silent clamp.
\item \texttt{StandardGenerator} is \emph{not} internally locked (deliberate); \texttt{HLCGenerator} \emph{is} internally locked (the \texttt{(l, c)} RMW must be atomic).
\item Randomness comes from the OS CSPRNG only: \texttt{getrandom(2)} / \texttt{getentropy(3)} / \texttt{BCryptGenRandom} / \texttt{os.urandom} / \texttt{SecureRandom} / \texttt{crypto.getRandomValues}. No userspace PRNG fallback.
\end{itemize}

\textbf{Conformance gate}: \texttt{tools/conformance/run.sh} is the branch-protection check on \texttt{main}. TesseraID, \tesserahash{}, and \tesseraseal{} all flow through the same harness; the latest cross-language run (2026-06-10) reports 0 failures across all eight languages.
